% ====================================================================
% FF3-5  히스 gap 닫힌꼴 ΔU^hys 의 거동 (Ch1 문맥, §4.2 제안 — §13 재참조 가능)
%   좌표 = eq:dUhys 그대로의 수치 평가 (gen_coords_ff3.py;
%   게이트: Ω=4RT@298.15K → 54.8 mV = fig:hysloop 캡션값 재현).
%   사용 라이브러리: arrows.meta, calc 만 (Ch1 preamble 내).
% ====================================================================
\begin{figure}[t]
\centering
\begin{tikzpicture}
% =============== (a) vs Omega at 298.15 K ===============
\begin{scope}[x=0.40cm,y=0.030cm]
\draw[-{Latex[length=1.5mm]}] (0,0) -- (0,132)
  node[above,font=\scriptsize] {$\Delta U_j^\hys$ [mV]};
\draw[-{Latex[length=1.5mm]}] (0,0) -- (15.3,0)
  node[below,font=\scriptsize] {$\Omega_j$ [kJ/mol]};
\foreach \xx in {5,10} {\draw (\xx,2.5) -- (\xx,-2.5) node[below,font=\scriptsize] {\xx};}
\foreach \yy in {40,80,120} {\draw (-0.15,\yy) -- (0.15,\yy) node[left,font=\scriptsize] {\yy};}
% threshold 2RT
\draw[densely dotted] (4.958,0) -- (4.958,128);
\node[font=\tiny,anchor=south,rotate=90] at (4.65,66) {$2RT\,(298\,\mathrm K)=4.96$};
% gap = 0 branch below threshold (exact zero)
\draw[very thick] (0,0) -- (4.958,0);
\node[font=\tiny,anchor=south] at (2.45,4) {gap $\equiv0$ ($\Omega_j\le2RT$)};
% closed-form curve above threshold
\draw[very thick] plot[smooth] coordinates {(4.958,0) (5.057,0.1926) (5.238,0.9066) (5.473,2.227) (5.751,4.191) (6.067,6.809) (6.415,10.08) (6.794,14) (7.202,18.56) (7.635,23.73) (8.094,29.51) (8.575,35.86) (9.08,42.78) (9.606,50.25) (10.15,58.24) (10.72,66.75) (11.3,75.75) (11.91,85.22) (12.53,95.17) (13.17,105.6) (13.83,116.4) (14,119.3)};
% staging initial values
\foreach \px/\py in {6.0/6.2, 8.0/28.3, 10.0/56.0, 13.0/102.8} {\fill (\px,\py) circle (1.1pt);}
\node[font=\tiny,anchor=west,align=left] at (9.6,20)
  {표~\ref{tab:staging} 초기값 $\Omega_j$\\$6/8/10/13$ kJ/mol\\$\to$ 상한 $6.2/28.3/56.0/102.8$ mV};
\draw[-{Latex[length=1.2mm]},black!55] (9.55,24) -- (8.15,28.0);
% onset annotation
\node[font=\tiny,anchor=west,align=left] at (5.35,105)
  {문턱 직상방 $\propto(\Omega-2RT)^{3/2}$\\--- $0$ 에서 \emph{양수로} 연속 발아};
\draw[-{Latex[length=1.2mm]},black!55] (5.9,96) -- (5.45,3.5);
\node[font=\scriptsize] at (2.2,120) {(a) $T=298.15$ K 고정};
\end{scope}
% =============== (b) vs T for three Omega ===============
\begin{scope}[xshift=4.45cm,x=0.0106cm,y=0.030cm]
\draw[-{Latex[length=1.5mm]}] (240,0) -- (240,132)
  node[above,font=\scriptsize] {$\Delta U_j^\hys$ [mV]};
\draw[-{Latex[length=1.5mm]}] (240,0) -- (820,0)
  node[below,font=\scriptsize] {$T$ [K]};
\foreach \xx in {300,450,600,750}
  {\draw (\xx,2.5) -- (\xx,-2.5) node[below,font=\scriptsize] {\xx};}
\foreach \yy in {40,80,120} {\draw (234,\yy) -- (246,\yy) node[left,font=\scriptsize] {\yy};}
\draw[densely dotted] (298.15,0) -- (298.15,116);
\node[font=\tiny,anchor=south] at (302,117) {$298$ K};
% Omega = 6000 (gray)
\draw[thick,black!55] plot[smooth] coordinates {(250,15.12) (255,14.05) (260.1,13.01) (265.1,12.01) (270.1,11.04) (275.2,10.09) (280.2,9.187) (285.3,8.312) (290.3,7.471) (295.3,6.663) (300.4,5.891) (305.4,5.154) (310.4,4.453) (315.5,3.791) (320.5,3.167) (325.6,2.585) (330.6,2.045) (335.6,1.551) (340.7,1.107) (345.7,0.7167) (350.7,0.389) (355.8,0.1371) (360.8,0)};
% Omega = 10000 (densely dashed)
\draw[thick,densely dashed] plot[smooth] coordinates {(250,71.7) (266,66.27) (281.9,61.06) (297.9,56.07) (313.9,51.28) (329.9,46.7) (345.8,42.32) (361.8,38.13) (377.8,34.13) (393.7,30.33) (409.7,26.71) (425.7,23.29) (441.7,20.05) (457.6,17.01) (473.6,14.17) (489.6,11.53) (505.5,9.091) (521.5,6.875) (537.5,4.891) (553.4,3.159) (569.4,1.71) (585.4,0.6013) (601.4,0)};
% Omega = 13000 (thick solid)
\draw[very thick] plot[smooth] coordinates {(250,121.3) (274.2,111.8) (298.3,102.7) (322.5,94.07) (346.7,85.84) (370.9,78) (395,70.54) (419.2,63.44) (443.4,56.69) (467.5,50.29) (491.7,44.22) (515.9,38.5) (540.1,33.1) (564.2,28.05) (588.4,23.33) (612.6,18.95) (636.7,14.93) (660.9,11.28) (685.1,8.018) (709.3,5.173) (733.4,2.798) (757.6,0.9828) (781.8,0)};
% T_c markers
\foreach \tc in {360.8,601.4,781.8} {\fill (\tc,0) circle (1.1pt);}
\node[font=\tiny,anchor=north] at (360.8,-4) {$T_c$};
\node[font=\tiny,anchor=north] at (601.4,-4) {$T_c$};
\node[font=\tiny,anchor=north] at (781.8,-4) {$T_c$};
% curve labels
\node[font=\tiny,anchor=west] at (403,74) {$\Omega{=}13$ kJ/mol};
\node[font=\tiny,anchor=west] at (382,27) {$10$};
\node[font=\tiny,anchor=west] at (295,13.5) {$6$};
\node[font=\tiny,align=left,anchor=west] at (500,108)
  {$T_c=\Omega_j/2R$ 에서\\$\propto(T_c-T)^{3/2}$ 소멸};
\node[font=\scriptsize] at (330,124) {(b) $\Omega_j$ 고정};
\end{scope}
\end{tikzpicture}
\caption{spinodal 상한 gap 닫힌꼴의 거동 --- 좌표는
식~\eqref{eq:dUhys} $\Delta U_j^\hys=\tfrac{2}{F}[\Omega_ju_j-2RT\,\mathrm{artanh}\,u_j]$,
$u_j=\sqrt{1-2RT/\Omega_j}$ 그대로의 수치 평가다. (a) $T{=}298.15$ K 고정: $\Omega_j\le2RT$ 면 gap 이
정확히 $0$ 인 명시 분기이고, 문턱을 넘으면 $\tfrac{8RT}{3F}u_j^3\propto(\Omega-2RT)^{3/2}$ 로 $0$
에서 양수로 연속 발아한다(두 급수를 함께 전개해야 얻는 부호 --- \S\ref{sec:hys} 의 Taylor 함정).
표~\ref{tab:staging} 초기값 네 점의 값($6.2$--$102.8$ mV)은 $\gamma_j{=}1$ 의 열역학 상한이며 실측
분기는 $\gamma_j$ 만큼 안쪽이다(식~\eqref{eq:Ubranch}). (b) $\Omega_j$ 고정: 온도가 오르면 gap 이
줄어 임계온도 $T_{c,j}=\Omega_j/2R$ 에서 $(T_c-T)^{3/2}$ 로 매끄럽게 소멸한다 --- $\Omega_j$ 를
낮추는 도핑이 같은 식 안에서 gap 을 닫는 경로(\S\ref{sec:lco-hys} 의 억제 극한)는 (a)에서 왼쪽으로,
(b)에서 곡선을 아래 곡선으로 옮기는 것과 같다.}
\label{fig:cand-ff3-5}
\end{figure}
